\documentclass{article}

% packages
\usepackage{amsmath, amsthm, thmtools, amsfonts, amssymb, luacode, catchfile, hyperref, ifthen, graphicx, bbm}
\ifcsname c@kobocompile\endcsname
	\usepackage[a5paper, total={1072pt, 1448pt}, margin=10pt, includeheadfoot]{geometry} % set page margins
\else
	\usepackage[a4paper, margin=50pt, includeheadfoot]{geometry}
\fi
\usepackage[shortlabels]{enumitem}
\usepackage[skip=3pt, indent=0pt]{parskip}

% language
\usepackage[bidi=basic, layout=tabular, provide=*]{babel}
\ifcsname c@english\endcsname
	\babelprovide[main, import]{english}
\else
	\babelprovide[main, import, alph=letters.gershayim]{hebrew}
	\setlist[enumerate,2]{label={\alph*.}}
	\babelprovide{rl}
\fi
\babelfont{rm}[Renderer=Harfbuzz]{Libertinus Serif}
\babelfont{sf}{Libertinus Sans}
\babelfont{tt}{Libertinus Mono}

% style
\ifcsname c@stoppagebreaks\endcsname
\else
	\AddToHook{cmd/section/before}{\clearpage}	% Add line break before section
\fi
\linespread{1.3}
\setcounter{secnumdepth}{0}		% Remove default number tags from sections, this won't do well with theorems
\AtBeginDocument{\setlength{\belowdisplayskip}{3pt}}
\AtBeginDocument{\setlength{\abovedisplayskip}{3pt}}
\graphicspath{ {../images/} }

% operators
\DeclareMathOperator\cis{cis}
\DeclareMathOperator\Sp{Span}
\DeclareMathOperator\tr{tr}
\DeclareMathOperator\im{Im}
\DeclareMathOperator\re{Re}
\DeclareMathOperator\diag{diag}
\DeclareMathOperator*\lowlim{\underline{lim}}
\DeclareMathOperator*\uplim{\overline{lim}}
\DeclareMathOperator\rng{rng}
\DeclareMathOperator\Sym{Sym}
\DeclareMathOperator\Arg{Arg}
\DeclareMathOperator\Log{Log}
\DeclareMathOperator\dom{dom}
\DeclareMathOperator\supp{Supp}
\DeclareMathOperator\var{Var}
\DeclareMathOperator\cov{Cov}
\DeclareMathOperator\End{End}
\DeclareMathOperator\id{id}
\DeclareMathOperator\cl{cl}
\DeclareMathOperator\gal{Gal}
\DeclareMathOperator\dist{dist}
\DeclareMathOperator\form{form}
\DeclareMathOperator\term{term}
\DeclareMathOperator\sent{sent}
\DeclareMathOperator\Th{Th}
\DeclareMathOperator\Hom{Hom}
\DeclareMathOperator\Aut{Aut}
\DeclareMathOperator\Char{char}
\DeclareMathOperator\vol{vol}
\DeclareMathOperator\GL{GL}
\DeclareMathOperator\SL{SL}
\DeclareMathOperator\lcm{lcm}

% commands
\newcommand{\cupdot}{\mathbin{\mathaccent\cdot\cup}}
\newcommand{\indicator}{\mathbbm{1}}
\newcommand{\AAA}[0]{\mathbb{A}}
\newcommand{\Aa}[0]{\mathcal{A}}
\newcommand{\BB}[0]{\mathbb{B}}
\newcommand{\Bb}[0]{\mathcal{B}}
\newcommand{\CC}[0]{\mathbb{C}}
\newcommand{\Cc}[0]{\mathcal{C}}
\newcommand{\DD}[0]{\mathbb{D}}
\newcommand{\Dd}[0]{\mathcal{D}}
\newcommand{\EE}[0]{\mathbb{E}}
\newcommand{\FF}[0]{\mathbb{F}}
\newcommand{\Ff}[0]{\mathcal{F}}
\newcommand{\Ii}[0]{\mathcal{I}}
\newcommand{\Gg}[0]{\mathcal{G}}
\newcommand{\HH}[0]{\mathbb{H}}
\newcommand{\Hh}[0]{\mathcal{H}}
\newcommand{\Kk}[0]{\mathcal{K}}
\newcommand{\Ll}[0]{\mathcal{L}}
\newcommand{\Mm}[0]{\mathcal{M}}
\newcommand{\NN}[0]{\mathbb{N}}
\newcommand{\Nn}[0]{\mathcal{N}}
\newcommand{\PP}[0]{\mathbb{P}}
\newcommand{\Pp}[0]{\mathcal{P}}
\newcommand{\QQ}[0]{\mathbb{Q}}
\newcommand{\RR}[0]{\mathbb{R}}
\newcommand{\Rr}[0]{\mathcal{R}}
\newcommand{\Ss}[0]{\mathcal{S}}
\newcommand{\TT}[0]{\mathbb{T}}
\newcommand{\Tt}[0]{\mathcal{T}}
\newcommand{\UU}[0]{\mathbb{U}}
\newcommand{\Uu}[0]{\mathcal{U}}
\newcommand{\Vv}[0]{\mathcal{V}}
\newcommand{\Ww}[0]{\mathcal{W}}
\newcommand{\Xx}[0]{\mathcal{X}}
\newcommand{\Yy}[0]{\mathcal{Y}}
\newcommand{\ZZ}[0]{\mathbb{Z}}
\newcommand{\Zz}[0]{\mathcal{Z}}
\newcommand{\acts}[0]{\circlearrowright}
\newcommand{\explain}[2] {
	\begin{flalign*}
		 && \text{#2} && \text{#1}
	\end{flalign*}
}
\newcommand{\maketitleprint}[1][green]{
	\begin{center}
		\includegraphics[width=6cm]{dragon_#1}
	\end{center}
}

% theorem commands
\newtheoremstyle{c_remark}
	{}	% Space above
	{}	% Space below
	{}% Body font
	{}	% Indent amount
	{\bfseries}	% Theorem head font
	{}	% Punctuation after theorem head
	{.5em}	% Space after theorem head
	{\thmname{#1}\thmnumber{ #2}\thmnote{ \normalfont{\text{(#3)}}}}	% head content
\newtheoremstyle{c_definition}
	{3pt}	% Space above
	{3pt}	% Space below
	{}% Body font
	{}	% Indent amount
	{\bfseries}	% Theorem head font
	{}	% Punctuation after theorem head
	{.5em}	% Space after theorem head
	{\thmname{#1}\thmnumber{ #2}\thmnote{ \normalfont{\text{(#3)}}}}	% head content
\newtheoremstyle{c_plain}
	{3pt}	% Space above
	{3pt}	% Space below
	{\itshape}% Body font
	{}	% Indent amount
	{\bfseries}	% Theorem head font
	{}	% Punctuation after theorem head
	{.5em}	% Space after theorem head
	{\thmname{#1}\thmnumber{ #2}\thmnote{ \text{(#3)}}}	% head content

\ifcsname c@english\endcsname
	\theoremstyle{plain}
	\newtheorem{theorem}{Theorem}[section]
	\newtheorem*{theorem*}{Theorem}
	\newtheorem{lemma}[theorem]{Lemma}
	\newtheorem{proposition}[theorem]{Proposition}
	\newtheorem*{proposition*}{Proposition}

	\theoremstyle{definition}
	\newtheorem{definition}[theorem]{Definition}
	\newtheorem*{definition*}{Definition}
	\newtheorem{example}{Example}[section]
	\newtheorem{exercise}{Exercise}[section]

	\theoremstyle{remark}
	\newtheorem*{remark}{Remark}
	\newtheorem*{solution}{Solution}
	\newtheorem{corollary}[theorem]{Corollary}
	\newtheorem{notation}[theorem]{Notation}
\else
	\theoremstyle{c_plain}
	\newtheorem{theorem}{משפט}[section]
	\newtheorem*{theorem*}{משפט}
	\newtheorem{lemma}[theorem]{למה}
	\newtheorem{proposition}[theorem]{טענה}
	\newtheorem*{proposition*}{טענה}

	\theoremstyle{c_definition}
	\newtheorem{definition}[theorem]{הגדרה}
	\newtheorem*{definition*}{הגדרה}
	\newtheorem{example}{דוגמה}[section]
	\newtheorem{exercise}{תרגיל}[section]

	\theoremstyle{c_remark}
	\newtheorem*{remark}{הערה}
	\newtheorem*{solution}{פתרון}
	\newtheorem{corollary}[theorem]{מסקנה}
	\newtheorem{notation}[theorem]{סימון}
\fi

% Questions related commands
\newcounter{question}
\setcounter{question}{1}
\newcounter{sub_question}
\setcounter{sub_question}{1}

\ifcsname c@english\endcsname
	\newcommand{\question}[1][0]{
		\ifthenelse{#1 = 0}{}{\setcounter{question}{#1}}
		\section{Question \arabic{question}}
		\addtocounter{question}{1}
		\setcounter{sub_question}{1}
	}

	\newcommand{\subquestion}[1][0]{
		\ifthenelse{#1 = 0}{}{\setcounter{sub_question}{#1}}
		\subsection{Part \alph{sub_question}}
		\addtocounter{sub_question}{1}
	}
\else
	\newcommand{\question}[1][0]{
		\ifthenelse{#1 = 0}{}{\setcounter{question}{#1}}
		\section{שאלה \arabic{question}}
		\addtocounter{question}{1}
		\setcounter{sub_question}{1}
	}

	\newcommand{\subquestion}[1][0]{
		\ifthenelse{#1 = 0}{}{\setcounter{sub_question}{#1}}
		\subsection{סעיף \localecounter{letters.gershayim}{sub_question}}
		\addtocounter{sub_question}{1}
	}
\fi

% import lua and start of document
\directlua{common = require ('../common')}

\GetEnv{AUTHOR}

% headers
\author{\AUTHOR}
\date\today

\title{פתרון מטלה 01 --- יסודות באלגברה קומוטטיבית}

\begin{document}
\maketitle
\maketitleprint[red]

\question{}
יהי $R$ חוג ויהי $0 < n \in \NN$, בשאלה זו נעסוק בחוג הפולינומים $R[x_1, \ldots, x_n]$, כאשר $\iota : R \to R[x_1, \ldots, x_n]$ היא השיכון.

\subquestion{}
נראה את התכונה האוניברסלית על חוג הפולינומים:
יהי $S$ חוג ויהי $\varphi : R \to S$ הומומורפיזם, נראה כי לכל $s_1, \ldots, s_n \in S$ קיים הומומורפיזם יחיד $\bar{\varphi} : R[x_1, \ldots, x_n] \to S$,
המקיים $\bar{\varphi} \circ \iota = \varphi$ וכן $\bar{\varphi}(x_i) = s_i$ לכל $i \in [n]$.
\begin{proof}
	נגדיר את $\bar{\varphi} : R[x_1, \ldots, x_n] \to S$ על־ידי,
	\[
		\forall i \in [n],\ \bar{\varphi}(x_i) = s_i,
		\qquad
		\forall r \in R,\ \bar{\varphi}(r) = \varphi(r)
	\]
	כאשר אנו מתייחסים ל־$r \in R$ גם כאל פולינום מדרגה $0$ בעל המקדם $r$ ב־$R[x_1, \ldots, x_n]$.
	ברור כי על־ידי סגירות לכפל וחיבור $\bar{\varphi}$ מוגדרת לכל פולינום ב־$R[x_1, \ldots, x_n]$, לכן זהו אכן הומומורפיזם חוגים, והוא מוגדר לקיים את התנאים הדרושים בדיוק.

	נבחין כי $R \cup \{ x_1, \ldots, x_n \}$ קבוצת יוצרים של חוג הפולינומים ולכן מספיק לקבוע את ערכם, ואכן מקביעת ערכים אלה נקבל גם יחידות של $\bar{\varphi}$.
\end{proof}

\subquestion{}
נראה כי מתקיים,
\[
	R[x_1, \ldots, x_{n - 1}][x_n]
	\cong R[x_1, \ldots, x_n]
\]
\begin{proof}
	נגדיר את $\varphi : R[x_1, \ldots, x_{n - 1}] \to R[x_1, \ldots, x_n]$ כשיכון, קרי $\varphi(x^i) = x^i$ לכל $i$. \\
	מסעיף א' נובע שקיים הומומורפיזם יחיד $\bar{\varphi} : R[x_1, \ldots, x_{n - 1}][x_n] \to R[x_1, \ldots, x_n]$ המקיים $\bar{\varphi} \circ \iota = \varphi$ וכן $\bar{\varphi}(x_i) = x_i$ לכל $i \in [n - 1]$.
	כלומר מתקיים $\bar{\varphi}(x_i) = x_i$ לכל $i \in [n]$ וכן $\bar{\varphi}(r) = r$ לכל $r \in R$.
	$\bar{\varphi}$ הוא על שכן לכל פולינום נוכל למצוא פולינום מתאים הנשלח אליו (ואשר בהצגתו בסמלים יהיה שקול לפולינום הרצוי ולכן לא נראה זאת פה).
	בנוסף $\bar{\varphi}$ חד־חד ערכי כהומומורפיזם שאיברי קבוצה יוצרת של מקורו נשלחים לערכים שונים, ולכן נסיק ש־$\bar{\varphi}$ איזומורפיזם.
\end{proof}

\question{}
יהי $R$ תחום שלמות, ויהי $0 < n \in \NN$. \\
נראה שלכל $f, g \in R[x_1, \ldots, x_n]$ מתקיים,
\[
	\deg fg
	= \deg f + \deg g
\]
\begin{proof}
	נגדיר יחס סדר $< \subseteq {\{ x_1^{r_1} \cdots x_n^{r_n} \mid r_1, \ldots, r_n \in \NN_0 \}}^2$ כך ש־$x_1^{r_1} \cdots x_n^{r_n} < x_1^{k_1} \cdots x_n^{k_n}$ אם,
	\[
		\sum_{i = 1}^n r_i
		< \sum_{i = 1}^n k_i
	\]
	או שהסכום שווה אבל קיים $i \in [n]$ כך ש־$r_i < r_k$ ולכל $j < i$ גם $r_j = k_j$, כלומר שהסדר של ה־$x_i$ המינימלי השונה בין המונומים הוא הקובע.

	נראה שזהו אכן יחס סדר: \\
	אנטי־רפלקסיביות: ברור כי $\sum_{i = 1}^n r_i \not< \sum_{i = 1}^n r_i$, וכן לא קיים אינדקס בו $r_i \ne r_i$. \\
	אסימטרי: נובע מאסימטריה של יחס הסדר על מספרים. \\
	טרנזיטיביות: נניח ש־$x_1^{r_1} \cdots x_n^{r_n} < x_1^{k_1} \cdots x_n^{k_n} < x_1^{l_1} \cdots x_n^{l_n}$, אז אם אחד מהיחסים האלה נקבע בשל סכומי סדרים שונים, אז סיימנו.
	נניח שסכום הדרגות שווה בכל שלושת המונומים.
	במקרה זה נסיק שקיים $i$ המקיים את התכונה שתיארנו עבור היחס בין שני המונומים הראשונים, $i$ זה בהכרח מקיים גם את היחס בין שני המונומים בקצוות.
	נסיק אם כך שאכן $<$ הוא יחס סדר חד.

	נניח שמתקיים,
	\[
		x_1^{r_1} \cdots x_n^{r_n} < x_1^{k_1} \cdots x_n^{k_n},
		\qquad
		x_1^{l_1} \cdots x_n^{l_n} < x_1^{m_1} \cdots x_n^{m_n}
	\]
	ונראה שגם,
	\[
		x_1^{r_1} \cdots x_n^{r_n} x_1^{l_1} \cdots x_n^{l_n}
		< x_1^{k_1} \cdots x_n^{k_n} x_1^{m_1} \cdots x_n^{m_n}
	\]
	כלומר, נראה שיחס הסדר נשמר תחת כפל.
	כמו בחלק הקודם, אם יחס הסדר נובע לפחות באחד הביטויים מיחס בסכום הסדרים, אז סיימנו, לכן נניח שסכומי החזקות זהים בשני המקרים.
	נבחר את $i_1, i_2$ המקיימים את התכונה השנייה, ונבחר $j = \min\{ i_1, i_2 \}$, אז הוא מקיים את התכונה הרצויה עבור המכפלות, זאת שכן $r_i l_i < k_i m_i$ מהגדרה.

	נניח ש־$f, g \in R[x_1, \ldots, x_n]$ ונניח ש־$x_1^{r_1} \cdots x_n^{r_n}$ ו־$x_1^{k_1} \cdots x_n^{k_n}$ המונומים הגדולים של $f, g$ בהתאמה, \\
	אז משימור יחס הסדר נקבל ש־$x_1^{r_1 + k_1} \cdots x_n^{r_n + k_n}$ הוא המונום המוביל של $f g$, ולכן נובע $\deg fg = \deg f + \deg g$ כפי שרצינו.
\end{proof}

\question{}
נראה שכל תחום שלמות סופי הוא שדה.
\begin{proof}
	יהי $(A, +, \cdot)$ תחום שלמות, כלומר חוג חילופי ללא מחלקי אפס, ונראה שהוא שדה. \\
	נבחין כי עלינו רק להראות ש־$(A^\times, \cdot)$ הוא חבורה כפלית, כל שאר התכונות של שדה ניתנות כבר מהגדרה של חוג חילופי.
	מההנחה ש־$A$ הוא תחום שלמות נובע ש־$A^\times = A \setminus \{ 0 \}$, וכן הנחנו שפעולת הכפל היא אסוציאטיבית וכי $1 \in A^\times$ הוא איבר נייטרלי, ולכן נשאר רק להראות שלכל איבר קיים הופכי כפלי.
	יהי $a \in A \setminus \{ 0 \}$, ונניח שקיימים גם $b, c \in A^\times$ כך ש־$a b = a c$, אז מתקיים $a (b - c) = 0$ ונקבל ש־$a$ הוא מחלק אפס בסתירה, לכן המיפוי $x \mapsto a x$ הוא חד־חד ערכי.
	מיפוי זה גם לא שולח לאפס, אחרת שוב $a$ יהיה מחלק אפס, לכן זהו אופרטור חד־חד ערכי ב־$A^\times$, ולכן מסופיות גם על, כלומר קיים $b \in A^\times$ כך ש־$a b = 1$ כפי שרצינו.
\end{proof}

\question{}
יהי $R$ חוג,
נראה ש־$P \triangleleft R$ הוא ראשוני אם ורק אם $R / P$ הוא תחום שלמות.
\begin{proof}
	$P$ ראשוני אם ורק אם $x y \in P \implies x \in P \lor y \in P$ אם ורק אם $x y + P = P \implies \bar{x} = P \lor \bar{y} = P$ אם ורק אם $x + P, y + P \ne P \implies x y + P \ne P$ וזוהי בדיוק הגדרת תחום שלמות $R / P$.
\end{proof}

\question{}
יהי $R$ חוג, ונאמר שאידאל $R \ne M \triangleleft R$ נקרא מקסימלי אם לכל אידאל $M \subsetneq I \triangleleft R$ מתקיים ש־$I = M$ או $I = R$. \\
נראה כי התנאים הבאים שקולים:
\begin{enumerate}
	\item $M \triangleleft R$ מקסימלי
	\item לכל $x \in R \setminus M$ מתקיים $(x) + M = R$
	\item $R / M$ שדה
\end{enumerate}
\begin{proof}
	$1 \implies 2$:
	נניח בשלילה שיש איבר שלא מקיים תכונה זו, כלומר $(x) + M \subsetneq R$, ונגדיר את האידאל $I = (\{ x \} \cup M)$.
	אז מהגדרה נקבל ש־$I \subsetneq R$ אבל $M \subsetneq I$ ונקבל סתירה למקסימליות $M$.

	$1 \implies 2$:
	זהה לחלק הקודם.

	$1 \implies 3$:
	\[
		\forall x \in R \setminus M,\ (x) + M = R
		\iff \forall x \in R \setminus M,\ 1 \in (x) + M
		\iff \forall x,\ (x + M) = R / M
	\]
	כלומר $M$ מקסימלי אם ורק אם כל אידאל ב־$R / M$ הוא ריק או כל $R / M$, ונותר להראות שזה נכון אם ורק אם זהו שדה, אך זוהי טענה ישירה משייכות יחידה לאידאלים לא ריקים.
\end{proof}

\question{}
\subquestion{}
יהי $n \in \NN$, נראה שמתקיים,
\[
	\ZZ[x] / (x^2 + 1, x + n)
	\cong \ZZ / (n^2 + 1)
\]
\begin{proof}
	נגדיר את המיפוי $\varphi : \ZZ[x] \to \ZZ$ על־ידי,
	\[
		\varphi(p)
		= p(-n)
	\]
	ברור כי זהו הומומורפיזם חוגים, וכן,
	\[
		\ker \varphi
		= \{ p \in \ZZ[x] \mid p(-n) = 0 \}
		= (x + n)
	\]
	וגם,
	\[
		\im \varphi
		= \{ p(-n) \mid p \in \ZZ[x] \}
		= \ZZ
	\]
	ולכן ממשפט האיזומורפיזם הראשון נובע $\ZZ[x] / (x + n) \cong \ZZ$.
	ממשפט האיזומורפיזם השלישי נובע,
	\[
		\ZZ / (x + n, x^2 + 1)
		\cong (\ZZ / (x + n)) / ((x + n, x^2 + 1) / (x + n))
	\]
	אבל,
	\[
		x^2 + 1 + (x + n)
		= n^2 + 1
	\]
	מבדיקה ישירה, ולכן סיימנו.
\end{proof}

\subquestion{}
נראה שאם $n^2 + 1$ ראשוני ב־$\ZZ$ אז $(n + i)$ אידאל מקסימלי ב־$\ZZ[i]$.
\begin{proof}
	נבחין כי $\ZZ[i] \cong \ZZ[x] / (x^2 + 1)$ ולכן מהסעיף הקודם ושאלה 4 הטענה נובעת.
\end{proof}

\question{}
יהיו $A, B$ חוגים ויהי $\varphi : A \to B$ הומומורפיזם.

\subquestion{}
נראה שאם $P \triangleleft B$ ראשוני אז גם $\varphi^{-1}(P) \triangleleft A$ ראשוני.
\begin{proof}
	ידוע לנו כבר כי $I = \varphi^{-1}(P)$ הוא אידאל ב־$A$, ונותר לנו רק להראות שהוא ראשוני.
	יהיו $x, y \in A$ כך ש־$x y \in I$ ונראה ש־$x \in I \lor y \in I$.
	נסמן $a = \varphi(x), b = \varphi(y)$ ולכן $a b \in P$ ומראשוניות $P$ נקבל בלי הגבלת הכלליות ש־$a \in P$ גם כן, אבל $I = \varphi^{-1}(P)$ ולכן בפרט $\varphi^{-1}(a) \subseteq I$ ובפרט $x \in I$.
\end{proof}

\subquestion{}
יהי $M \triangleleft B$ אידאל מקסימלי, ונראה ש־$A = \varphi^{-1}(M) \triangleleft A$ הוא לא בהכרח מקסימלי.
\begin{solution}
	נגדיר $A = \ZZ, B = \ZZ[x]$ וכן $P = (x)$, אז $P \triangleleft B$ הוא מקסימלי כטענה מהתרגול.
	נגדיר גם $\varphi = i$ השיכון הטריוויאלי, לכן $\varphi^{-1}(P) = \{ 0 \}$ בלבד, וזהו כמובן אידאל לא מקסימלי ב־$A$.
\end{solution}

\subquestion{}
נניח ש־$\varphi$ הוא על ונראה שאם $M \triangleleft B$ מקסימלי אז גם $\varphi^{-1}(M)$ מקסימלי.
\begin{proof}
	נסמן שוב $I = \varphi^{-1}(M)$ ויהי $I \subsetneq J \subset A$ אידאל נוסף.
	מהטענה כי אם $\varphi$ על אז הוא משמר אידאלים בתמונה נסיק ש־$J' = \varphi(J)$ הוא אידאל ב־$B$, אבל $\varphi(I) \subsetneq \varphi(J)$, כלומר $M \subsetneq J' \subset B$ ולכן ממקסימליות $M$ נקבל ש־$J' = B$ ולכן $J = A$.
\end{proof}

\question{}
יהי $R$ חוג ו־$I \triangleleft R$.
נראה ש־$R[x] / I[x] \cong (R / I)[x]$ ונסיק שאם $P \triangleleft R$ ראשוני אז גם $P[x] \triangleleft R[x]$ ראשוני.
\begin{proof}
	נגדיר את המיפוי,
	\[
		\varphi : R[x] / I[x] \to (R / I)[x],
		\qquad
		\varphi(a_0 x^0 + \cdots + x_n x^n + I[x])
		= (a_0 + I) x^0 + \cdots + (a_n + I) x^n
	\]
	ונראה שהוא מוגדר היטב ואיזומורפיזם. \\
	ברור כי הוא מוגדר היטב על־ידי שימוש בהגדרת $I[x]$ ונעבור לבדוק שזהו הומומורפיזם הפיך. \\
	על נובע ישירות מהגדרה, עבור בחירת נציגים בתמונה ניתן להשתמש בהם בתחום, חד־חד ערכיות נובעת אף היא מאותה סיבה, לכן זהו איזומורפיזם וסיימנו.

	נניח ש־$P \triangleleft R$ הוא ראשוני,
	משאלה 4 נובע ש־$R / P$ הוא תחום שלמות, ולכן גם $(R / P)[x]$ תחום שלמות, ותוך שימוש בזהות שמצאנו זה עתה נסיק שגם $R[x] / P[x]$ תחום שלמות ולכן $P[x] \triangleleft R[x]$ ראשוני.
\end{proof}

\question{}
יהי $R$ חוג כך שלכל $x \in R$ קיים $1 < n \in \NN$ כך ש־$x^n = x$. \\
נראה שכל אידאל ראשוני של $R$ הוא מקסימלי.
\begin{proof}
	יהי אידאל ראשוני $P \triangleleft R$ ויהי $x \in R \setminus P$.
	נסמן $K = R / P$ ונסיק שזהו תחום שלמות, נראה שלכל $\bar{x} \in K^\times$ יש הופכי כפלי.
	קיים $n$ עבורו $x^n = x$ ולכן גם $\bar{x}^n = \bar{x}$.
	נסמן $y = x^{n - 1}$ ונקבל ש־$y x = x \iff x (y - 1) = 0$, אז $\bar{z} = (y - 1) + P$ הוא איבר המקיים $\bar{x} \bar{z} = 0$, אבל $K$ תחום שלמות ולכן $\bar{x} = 0$ או $\bar{z} = 0$, ולכן בהכרח $\bar{y} = \bar{1}$ בלבד.
	נובע אם כך שיש ל־$\bar{x}$ הופכי, לכן $K$ שדה ונובע ש־$P$ אידאל מקסימלי.
\end{proof}

\question{}
יהי $R$ חוג כך ש־$x^2 = x$ לכל $x \in R$.

\subquestion{}
נראה שלכל $x \in R$ מתקיים $2 x = 0$.
\begin{proof}
	נשתמש בהנחה שלנו על $x + 1, x$ ונקבל,
	\[
		x + 1
		= {(x + 1)}^2
		= x^2 + 2x + 1
		= x + 2 x + 1
	\]
	ולאחר העברת אגפים $2 x = 0$ בדיוק.
\end{proof}

\subquestion{}
נראה שלכל אידאל ראשוני $P \triangleleft R$ השדה $R / P$ הוא בעל שני איברים בדיוק.
\begin{proof}
	ידוע לנו כי זהו שדה מהשאלה הקודמת.
	ידוע כי $0, 1 \in R / P$, ולכל $x \in R / P$ מתקיים $x (x - 1) = 0$ ולכן $x \in \{0, 1\}$ וסיימנו.
\end{proof}

\subquestion{}
נראה שלכל $S \subseteq R$ סופי, $(S)$ הוא אידאל ראשי.
\begin{proof}
	מספיק להראות שאם $x, y \in R$ אז $(x, y) = (z)$ לאיזשהו $z \in R$.
	נגדיר $z = x + y + x y$ ולכן ברור כי $(z) \supset (x, y)$, ונראה גם את הכיוון ההפוך.
	\[
		z x
		= x^2 + x y + x y
		= x
	\]
	ולכן $x \in (z)$, ובאותו אופן נקבל שגם $y \in (z)$.
\end{proof}

\question{}
יהי $d \in \ZZ$ חסר ריבועים.
נראה ש־$a + b \sqrt{d} \in \ZZ[\sqrt{d}]$ הפיך אם ורק אם $a^2 - d b^2 \in \{ \pm 1 \}$.
\begin{proof}
	נניח ש־$a + b \sqrt{d}$ הפיך ונניח שההופכי שלו נתון על־ידי $x + y \sqrt{d}$, כלומר,
	\[
		1
		= (a + b \sqrt{d})(x + y \sqrt{d})
		= a x + b y d + \sqrt{d} (a y + b x)
	\]
	ולכן $a x = 1 - b y d$ וכן $a y = - b x$.
	\[
		a^2 x
		= a - a b y d
		= a - b d (-b x)
		= a + b^2 d x
		\implies x (a^2 - b^2 d)
		= a
	\]

	בכיוון ההפוך נבחין כי מתקיים,
	\[
		(a + b \sqrt{d})(a - b \sqrt{d})
		= a^2 - b^2 d
		= \pm 1
	\]
	ולכן $a - b \sqrt{d}$ או $-a + b \sqrt{d}$ ההופכי של $a + b \sqrt{d}$.

	דרך נוספת: \\
	נבחין כי $\ZZ[\sqrt{d}] \subset \RR$ תת־חוג, בפרט אם ל־$a + b \sqrt{d} \in \ZZ[d]$ יש הופכי אם ורק אם קיים הופכי $\in \RR$ כך שהוא שייך ל־$\ZZ[\sqrt{d}]$.
	מתקיים,
	\[
		r (a + b \sqrt{d}) = 1
		\iff
		r = \frac{a - b \sqrt{d}}{a^2 - b^2 d}
	\]
	כלומר $r \in \ZZ[\sqrt{d}] \iff a^2 - b^2 d \in \{ \pm 1 \}$.
\end{proof}

\end{document}
